% Source PDF: source/普通高考/2015/2015湖南理.pdf, page 1, question 3.
% The source flowchart is vector page content; original PNG remains at img/2015/hunan_science/q3_fig1.png.
% Node -> directed-edge list: 开始 -> 输入 n -> i=1,S=0 ->
% S=S+1/((2i-1)(2i+1)) -> i=i+1 -> i>n?;
% 是 -> 输出 S -> 结束; 否 -> right-side return route -> the stem immediately above the update box.
% All connectors are solid; no dashed segments or other edges. Node text is centered with local parindent=0pt.

\begin{tikzpicture}[
  >=Stealth,
  x=1cm, y=1cm,
  line width=0.45pt,
  line cap=round,
  line join=round,
  font=\normalsize,
  flowtext/.style={anchor=center,align=center,inner sep=0pt,
    execute at begin node={\setlength{\parindent}{0pt}}},
  flowbox/.style={flowtext,draw,minimum height=.68cm,inner xsep=3pt,inner ysep=2pt},
  terminal/.style={flowbox,rounded corners=2pt,minimum width=1.35cm},
  process/.style={flowbox},
  io/.style={flowbox,shape=trapezium,trapezium left angle=72,trapezium right angle=108,
    minimum height=.76cm},
  decision/.style={flowbox,shape=diamond,aspect=2.8,minimum width=3.20cm,
    minimum height=1.02cm,inner sep=0pt},
  branchlabel/.style={flowtext,inner sep=0pt}
]
  \node[terminal] (start) at (0,12.20) {开始};
  \node[io,minimum width=1.80cm] (input) at (0,10.90) {输入 $n$};
  \node[process,minimum width=2.40cm] (init) at (0,9.55) {$i=1,S=0$};
  \node[process,minimum width=4.60cm,minimum height=1.02cm] (update) at (0,7.75)
    {$S=S+\dfrac{1}{(2i-1)(2i+1)}$};
  \node[process,minimum width=1.85cm] (increment) at (0,6.20) {$i=i+1$};
  \node[decision] (test) at (0,4.75) {$i>n?$};
  \node[io,minimum width=2.15cm] (output) at (0,3.25) {输出 $S$};
  \node[terminal] (finish) at (0,1.85) {结束};

  \draw[->] (start.south) -- (input.north);
  \draw[->] (input.south) -- (init.north);
  \draw[->] (init.south) -- (update.north);
  \draw[->] (update.south) -- (increment.north);
  \draw[->] (increment.south) -- (test.north);
  \draw[->] (test.south) -- node[branchlabel,right] {是} (output.north);
  \draw[->] (output.south) -- (finish.north);

  % 否 returns independently along the right side and re-enters above update.
  \draw (test.east) -- node[branchlabel,above] {否} (3.20,4.75) -- (3.20,8.55);
  \draw[->] (3.20,8.55) -- (0,8.55);

  \path[use as bounding box] (-0.35,1.50) rectangle (3.55,12.55);
\end{tikzpicture}
